% ID: SEL-D0072-OA-E02
% ESTADO: validada. S072.
\textbf{a)} Por la regla del cociente, $f'(x)=1/(2x-1)^2$. En $(0,1)$ la pendiente es $1$, así que la tangente es $\boxed{y=x+1}$.

\textbf{b)} La derivada es positiva en su dominio $\mathbb R\setminus\{1/2\}$. La función crece estrictamente en cada intervalo $(-\infty,1/2)$ y $(1/2,\infty)$.

\textbf{c)} Escribimos $f(x)=1/2-1/[4(x-1/2)]$. En $1/2$ el límite izquierdo es $+\infty$ y el derecho es $-\infty$. En ambos infinitos el límite es $1/2$. Las asíntotas son $x=1/2$ e $y=1/2$. El numerador se anula en $1$, por lo que los cortes son $(1,0)$ y $(0,1)$. La rama izquierda queda por encima de la asíntota horizontal y la derecha por debajo.
\begin{center}
\setlength{\unitlength}{0.9mm}
\begin{picture}(114,68)(-48,-30)
\qbezier(-44,0)(7,0)(58,0)
\qbezier(0,-27)(0,1)(0,29)
{\color{gray}\qbezier(10,-27)(10,1)(10,29)
\qbezier(-44,5)(7,5)(58,5)}
\qbezier(-42,5.96154)(-35,6.08356)(-28,6.31579)
\qbezier(-28,6.31579)(-21,6.52625)(-14,7.08333)
\qbezier(-14,7.08333)(-7,7.34069)(0,10)
\qbezier(0,10)(2,10.8333)(4,13.3333)
\qbezier(4,13.3333)(5.2,14.7222)(6.4,18.8889)
\qbezier(6.4,18.8889)(7,20.9722)(7.6,25.8333)
\qbezier(11.6,-26.25)(12.1,-17.3787)(12.6,-14.2308)
\qbezier(12.6,-14.2308)(13.6,-7.72761)(14.6,-5.86957)
\qbezier(14.6,-5.86957)(17.3,-0.763848)(20,0)
\qbezier(20,0)(26,2.38636)(32,2.72727)
\qbezier(32,2.72727)(43.5,3.70684)(55,3.88889)
\put(59,0){$x$}\put(1,29){$y$}
\put(12,27){\scriptsize $x=1/2$}
\put(48,9){\scriptsize $y=1/2$}
\put(-12,12){\scriptsize $(0,1)$}
\put(22,-6){\scriptsize $(1,0)$}
\end{picture}
\end{center}
